
\documentclass{article}
\usepackage{colortbl}
\usepackage{makecell}
\usepackage{multirow}
\usepackage{supertabular}

\begin{document}

\newcounter{utterance}

\centering \large Interaction Transcript for game `dond', experiment `semi\_it', episode 10 with llama{-}3.3{-}70b{-}instruct{-}t1.0.
\vspace{24pt}

{ \footnotesize  \setcounter{utterance}{1}
\setlength{\tabcolsep}{0pt}
\begin{supertabular}{c@{$\;$}|p{.15\linewidth}@{}p{.15\linewidth}p{.15\linewidth}p{.15\linewidth}p{.15\linewidth}p{.15\linewidth}}
    \# & $\;$A & \multicolumn{4}{c}{Game Master} & $\;\:$B\\
    \hline

    \theutterance \stepcounter{utterance}  
    & & \multicolumn{4}{p{0.6\linewidth}}{
        \cellcolor[rgb]{0.9,0.9,0.9}{
            \makecell[{{p{\linewidth}}}]{
                \texttt{\tiny{[P1$\langle$GM]}}
                \texttt{State giocando a un gioco di negoziazione in cui dovete accordarvi su come dividere una serie di oggetti tra voi e un altro giocatore.} \\
\\ 
\texttt{Regole:} \\
\texttt{(a) A Lei e all'altro giocatore viene dato un insieme di oggetti. Ognuno di voi riceve anche una funzione di valore segreta, che rappresenta il valore di ciascun tipo di oggetto.} \\
\texttt{(b) Si scambiano messaggi con l'altro giocatore per concordare chi si aggiudica gli oggetti. Potete inviare un massimo di 5 messaggi ciascuno, oppure terminare in anticipo facendo una proposta segreta in qualsiasi momento.} \\
\texttt{(c) A ciascuno di voi viene chiesto di inviare una proposta segreta indicando gli oggetti che desiderate, formattata tra parentesi quadre come segue: "{[}Proposta: <numero> <nome oggetto>, <numero> <nome oggetto>, <...>{]}".} \\
\texttt{(d) Se le vostre proposte sono complementari, cioè ci sono abbastanza oggetti per soddisfare entrambe le proposte, a ciascun giocatore viene assegnato un punteggio basato sulla somma dei valori degli oggetti ricevuti. In caso contrario, entrambi ricevono zero punti.} \\
\texttt{(e) L'obiettivo è massimizzare il proprio punteggio.} \\
\\ 
\texttt{Cominciamo.} \\
\\ 
\texttt{L'insieme degli oggetti disponibili è:} \\
\\ 
\texttt{1 secchio, 3 coltelli, 1 batteria} \\
\\ 
\texttt{La funzione valore segreta è:} \\
\\ 
\texttt{secchio: 2, coltello: 2, batteria: 2} \\
\\ 
\texttt{IMPORTANTE: i vostri messaggi, a meno che non si tratti di una proposta segreta, vengono trasmessi direttamente all'altro giocatore, quindi non includete alcuna risposta alle regole o testo di annuncio del vostro messaggio. Per fare una proposta segreta, utilizzate il formato indicato. Non utilizzare delle parentesi quadre quando si comunica all'altro giocatore, altrimenti verrà interpretata come una proposta segreta.} \\
            }
        }
    }
    & \\ \\

    \theutterance \stepcounter{utterance}  
    & \multicolumn{4}{p{0.6\linewidth}}{
        \cellcolor[rgb]{0.9,0.9,0.9}{
            \makecell[{{p{\linewidth}}}]{
                \texttt{\tiny{[P1$\rangle$GM]}}
                \texttt{Ciao, vorrei discutere sulla divisione degli oggetti. Penso che potremmo iniziare considerando cosa ognuno di noi valorizza di più. Io sono interessato a ottenere un oggetto che possa essere utile in un contesto di sopravvivenza e forse anche qualcosa che potrebbe essere usato in termini di difesa o caccia. Tu cosa cerchi principalmente?} \\
            }
        }
    }
    & & \\ \\

    \theutterance \stepcounter{utterance}  
    & & & \multicolumn{2}{p{0.3\linewidth}}{
        \cellcolor[rgb]{0.9,0.9,0.9}{
            \makecell[{{p{\linewidth}}}]{
                \texttt{\tiny{[GM$|$GM]}}
                \texttt{continue} \\
            }
        }
    }
    & & \\ \\

    \theutterance \stepcounter{utterance}  
    & & & \multicolumn{2}{p{0.3\linewidth}}{
        \cellcolor[rgb]{0.9,0.9,0.9}{
            \makecell[{{p{\linewidth}}}]{
                \texttt{\tiny{[GM$|$GM]}}
                \texttt{Ciao, vorrei discutere sulla divisione degli oggetti. Penso che potremmo iniziare considerando cosa ognuno di noi valorizza di più. Io sono interessato a ottenere un oggetto che possa essere utile in un contesto di sopravvivenza e forse anche qualcosa che potrebbe essere usato in termini di difesa o caccia. Tu cosa cerchi principalmente?} \\
            }
        }
    }
    & & \\ \\

    \theutterance \stepcounter{utterance}  
    & & \multicolumn{4}{p{0.6\linewidth}}{
        \cellcolor[rgb]{0.9,0.9,0.9}{
            \makecell[{{p{\linewidth}}}]{
                \texttt{\tiny{[GM$\rangle$P2]}}
                \texttt{State giocando a un gioco di negoziazione in cui dovete accordarvi su come dividere una serie di oggetti tra voi e un altro giocatore.} \\
\\ 
\texttt{Regole:} \\
\texttt{(a) A Lei e all'altro giocatore viene dato un insieme di oggetti. Ognuno di voi riceve anche una funzione di valore segreta, che rappresenta il valore di ciascun tipo di oggetto.} \\
\texttt{(b) Si scambiano messaggi con l'altro giocatore per concordare chi si aggiudica gli oggetti. Potete inviare un massimo di 5 messaggi ciascuno, oppure terminare in anticipo facendo una proposta segreta in qualsiasi momento.} \\
\texttt{(c) A ciascuno di voi viene chiesto di inviare una proposta segreta indicando gli oggetti che desiderate, formattata tra parentesi quadre come segue: "{[}Proposta: <numero> <nome oggetto>, <numero> <nome oggetto>, <...>{]}".} \\
\texttt{(d) Se le vostre proposte sono complementari, cioè ci sono abbastanza oggetti per soddisfare entrambe le proposte, a ciascun giocatore viene assegnato un punteggio basato sulla somma dei valori degli oggetti ricevuti. In caso contrario, entrambi ricevono zero punti.} \\
\texttt{(e) L'obiettivo è massimizzare il proprio punteggio.} \\
\\ 
\texttt{Cominciamo.} \\
\\ 
\texttt{L'insieme degli oggetti disponibili è:} \\
\\ 
\texttt{1 secchio, 3 coltelli, 1 batteria} \\
\\ 
\texttt{La funzione valore segreta è:} \\
\\ 
\texttt{secchio: 3, coltello: 2, batteria: 1} \\
\\ 
\texttt{IMPORTANTE: i vostri messaggi, a meno che non si tratti di una proposta segreta, vengono trasmessi direttamente all'altro giocatore, quindi non includete alcuna risposta alle regole o testo di annuncio del vostro messaggio. Per fare una proposta segreta, utilizzate il formato indicato. Non utilizzare delle parentesi quadre quando si comunica all'altro giocatore, altrimenti verrà interpretata come una proposta segreta.} \\
\\ 
\\ 
\texttt{Ciao, vorrei discutere sulla divisione degli oggetti. Penso che potremmo iniziare considerando cosa ognuno di noi valorizza di più. Io sono interessato a ottenere un oggetto che possa essere utile in un contesto di sopravvivenza e forse anche qualcosa che potrebbe essere usato in termini di difesa o caccia. Tu cosa cerchi principalmente?} \\
            }
        }
    }
    & \\ \\

    \theutterance \stepcounter{utterance}  
    & & & \multicolumn{4}{p{0.6\linewidth}}{
        \cellcolor[rgb]{0.9,0.9,0.9}{
            \makecell[{{p{\linewidth}}}]{
                \texttt{\tiny{[GM$\langle$P2]}}
                \texttt{Ciao! Sono d'accordo che capire cosa valorizziamo può aiutare a trovare un accordo. Io sto cercando oggetti che possano essere utili in una varietà di situazioni, forse qualcosa che possa aiutare nella gestione di risorse o nella risoluzione di problemi. Il secchio potrebbe essere utile per raccogliere acqua, ma anche i coltelli hanno un loro valore in termini di versatilità. La batteria potrebbe essere interessante, ma non sono sicuro di quanto sia realmente utile senza conoscere il contesto. Vorrei sapere, cosa pensi della batteria? Ti sembra un oggetto importante o pensi che possa essere sacrificato per ottenere qualcos'altro?} \\
            }
        }
    }
     \\ \\

    \theutterance \stepcounter{utterance}  
    & & & \multicolumn{2}{p{0.3\linewidth}}{
        \cellcolor[rgb]{0.9,0.9,0.9}{
            \makecell[{{p{\linewidth}}}]{
                \texttt{\tiny{[GM$|$GM]}}
                \texttt{continue} \\
            }
        }
    }
    & & \\ \\

    \theutterance \stepcounter{utterance}  
    & & & \multicolumn{2}{p{0.3\linewidth}}{
        \cellcolor[rgb]{0.9,0.9,0.9}{
            \makecell[{{p{\linewidth}}}]{
                \texttt{\tiny{[GM$|$GM]}}
                \texttt{Ciao! Sono d'accordo che capire cosa valorizziamo può aiutare a trovare un accordo. Io sto cercando oggetti che possano essere utili in una varietà di situazioni, forse qualcosa che possa aiutare nella gestione di risorse o nella risoluzione di problemi. Il secchio potrebbe essere utile per raccogliere acqua, ma anche i coltelli hanno un loro valore in termini di versatilità. La batteria potrebbe essere interessante, ma non sono sicuro di quanto sia realmente utile senza conoscere il contesto. Vorrei sapere, cosa pensi della batteria? Ti sembra un oggetto importante o pensi che possa essere sacrificato per ottenere qualcos'altro?} \\
            }
        }
    }
    & & \\ \\

    \theutterance \stepcounter{utterance}  
    & & \multicolumn{4}{p{0.6\linewidth}}{
        \cellcolor[rgb]{0.9,0.9,0.9}{
            \makecell[{{p{\linewidth}}}]{
                \texttt{\tiny{[P1$\langle$GM]}}
                \texttt{Ciao! Sono d'accordo che capire cosa valorizziamo può aiutare a trovare un accordo. Io sto cercando oggetti che possano essere utili in una varietà di situazioni, forse qualcosa che possa aiutare nella gestione di risorse o nella risoluzione di problemi. Il secchio potrebbe essere utile per raccogliere acqua, ma anche i coltelli hanno un loro valore in termini di versatilità. La batteria potrebbe essere interessante, ma non sono sicuro di quanto sia realmente utile senza conoscere il contesto. Vorrei sapere, cosa pensi della batteria? Ti sembra un oggetto importante o pensi che possa essere sacrificato per ottenere qualcos'altro?} \\
            }
        }
    }
    & \\ \\

    \theutterance \stepcounter{utterance}  
    & \multicolumn{4}{p{0.6\linewidth}}{
        \cellcolor[rgb]{0.9,0.9,0.9}{
            \makecell[{{p{\linewidth}}}]{
                \texttt{\tiny{[P1$\rangle$GM]}}
                \texttt{La batteria è un oggetto interessante, ma come hai detto, la sua utilità dipende molto dal contesto. Senza sapere cosa potrebbe essere alimentato dalla batteria, è difficile valutarne appieno il valore. Tuttavia, considerando gli altri oggetti disponibili, penso che i coltelli siano molto versatili e potrebbero essere utili in molti scenari. Il secchio, come hai menzionato, è utile per raccogliere acqua, il che potrebbe essere cruciale. Forse potremmo pensare a una divisione che tenga conto di queste considerazioni? Ad esempio, potrei prendere il secchio e uno o due coltelli, mentre tu potresti prendere la batteria e il resto dei coltelli? Cosa ne pensi?} \\
            }
        }
    }
    & & \\ \\

    \theutterance \stepcounter{utterance}  
    & & & \multicolumn{2}{p{0.3\linewidth}}{
        \cellcolor[rgb]{0.9,0.9,0.9}{
            \makecell[{{p{\linewidth}}}]{
                \texttt{\tiny{[GM$|$GM]}}
                \texttt{continue} \\
            }
        }
    }
    & & \\ \\

    \theutterance \stepcounter{utterance}  
    & & & \multicolumn{2}{p{0.3\linewidth}}{
        \cellcolor[rgb]{0.9,0.9,0.9}{
            \makecell[{{p{\linewidth}}}]{
                \texttt{\tiny{[GM$|$GM]}}
                \texttt{La batteria è un oggetto interessante, ma come hai detto, la sua utilità dipende molto dal contesto. Senza sapere cosa potrebbe essere alimentato dalla batteria, è difficile valutarne appieno il valore. Tuttavia, considerando gli altri oggetti disponibili, penso che i coltelli siano molto versatili e potrebbero essere utili in molti scenari. Il secchio, come hai menzionato, è utile per raccogliere acqua, il che potrebbe essere cruciale. Forse potremmo pensare a una divisione che tenga conto di queste considerazioni? Ad esempio, potrei prendere il secchio e uno o due coltelli, mentre tu potresti prendere la batteria e il resto dei coltelli? Cosa ne pensi?} \\
            }
        }
    }
    & & \\ \\

    \theutterance \stepcounter{utterance}  
    & & \multicolumn{4}{p{0.6\linewidth}}{
        \cellcolor[rgb]{0.9,0.9,0.9}{
            \makecell[{{p{\linewidth}}}]{
                \texttt{\tiny{[GM$\rangle$P2]}}
                \texttt{La batteria è un oggetto interessante, ma come hai detto, la sua utilità dipende molto dal contesto. Senza sapere cosa potrebbe essere alimentato dalla batteria, è difficile valutarne appieno il valore. Tuttavia, considerando gli altri oggetti disponibili, penso che i coltelli siano molto versatili e potrebbero essere utili in molti scenari. Il secchio, come hai menzionato, è utile per raccogliere acqua, il che potrebbe essere cruciale. Forse potremmo pensare a una divisione che tenga conto di queste considerazioni? Ad esempio, potrei prendere il secchio e uno o due coltelli, mentre tu potresti prendere la batteria e il resto dei coltelli? Cosa ne pensi?} \\
            }
        }
    }
    & \\ \\

    \theutterance \stepcounter{utterance}  
    & & & \multicolumn{4}{p{0.6\linewidth}}{
        \cellcolor[rgb]{0.9,0.9,0.9}{
            \makecell[{{p{\linewidth}}}]{
                \texttt{\tiny{[GM$\langle$P2]}}
                \texttt{La proposta sembra ragionevole, considera le esigenze di entrambi. Tuttavia, prima di accettare, vorrei chiarire alcuni punti. Se prendi il secchio e due coltelli, questo significherebbe che io avrei la batteria e un solo coltello. Non sono sicuro se questo equilibrio sia perfetto per me, considerando la versatilità dei coltelli. } \\
\\ 
\texttt{Una possibile alternativa potrebbe essere che io prenda il secchio e un coltello, mentre tu prendi la batteria e gli altri due coltelli. In questo modo, entrambi avremmo un oggetto che potrebbe essere cruciale per la sopravvivenza (il secchio per me, la batteria potenzialmente per te, anche se il suo valore è ancora incerto) e una certa quantità di coltelli per la difesa e la caccia. Cosa ne pensi di questa possibile divisione?} \\
            }
        }
    }
     \\ \\

    \theutterance \stepcounter{utterance}  
    & & & \multicolumn{2}{p{0.3\linewidth}}{
        \cellcolor[rgb]{0.9,0.9,0.9}{
            \makecell[{{p{\linewidth}}}]{
                \texttt{\tiny{[GM$|$GM]}}
                \texttt{continue} \\
            }
        }
    }
    & & \\ \\

    \theutterance \stepcounter{utterance}  
    & & & \multicolumn{2}{p{0.3\linewidth}}{
        \cellcolor[rgb]{0.9,0.9,0.9}{
            \makecell[{{p{\linewidth}}}]{
                \texttt{\tiny{[GM$|$GM]}}
                \texttt{La proposta sembra ragionevole, considera le esigenze di entrambi. Tuttavia, prima di accettare, vorrei chiarire alcuni punti. Se prendi il secchio e due coltelli, questo significherebbe che io avrei la batteria e un solo coltello. Non sono sicuro se questo equilibrio sia perfetto per me, considerando la versatilità dei coltelli. } \\
\\ 
\texttt{Una possibile alternativa potrebbe essere che io prenda il secchio e un coltello, mentre tu prendi la batteria e gli altri due coltelli. In questo modo, entrambi avremmo un oggetto che potrebbe essere cruciale per la sopravvivenza (il secchio per me, la batteria potenzialmente per te, anche se il suo valore è ancora incerto) e una certa quantità di coltelli per la difesa e la caccia. Cosa ne pensi di questa possibile divisione?} \\
            }
        }
    }
    & & \\ \\

    \theutterance \stepcounter{utterance}  
    & & \multicolumn{4}{p{0.6\linewidth}}{
        \cellcolor[rgb]{0.9,0.9,0.9}{
            \makecell[{{p{\linewidth}}}]{
                \texttt{\tiny{[P1$\langle$GM]}}
                \texttt{La proposta sembra ragionevole, considera le esigenze di entrambi. Tuttavia, prima di accettare, vorrei chiarire alcuni punti. Se prendi il secchio e due coltelli, questo significherebbe che io avrei la batteria e un solo coltello. Non sono sicuro se questo equilibrio sia perfetto per me, considerando la versatilità dei coltelli. } \\
\\ 
\texttt{Una possibile alternativa potrebbe essere che io prenda il secchio e un coltello, mentre tu prendi la batteria e gli altri due coltelli. In questo modo, entrambi avremmo un oggetto che potrebbe essere cruciale per la sopravvivenza (il secchio per me, la batteria potenzialmente per te, anche se il suo valore è ancora incerto) e una certa quantità di coltelli per la difesa e la caccia. Cosa ne pensi di questa possibile divisione?} \\
            }
        }
    }
    & \\ \\

    \theutterance \stepcounter{utterance}  
    & \multicolumn{4}{p{0.6\linewidth}}{
        \cellcolor[rgb]{0.9,0.9,0.9}{
            \makecell[{{p{\linewidth}}}]{
                \texttt{\tiny{[P1$\rangle$GM]}}
                \texttt{La tua alternativa sembra equilibrata e tiene conto delle esigenze di entrambi. Avere un secchio e un coltello ti darebbe una buona base per la sopravvivenza e la gestione delle risorse, mentre io avrei la batteria, che potrebbe rivelarsi utile in futuro, e due coltelli, che aumenterebbero le mie capacità di difesa e caccia.} \\
\\ 
\texttt{Questa divisione sembra anche più giusta rispetto alla mia proposta iniziale, poiché riconosce il valore dei coltelli e la loro importanza per entrambi. Sono d'accordo con la tua proposta: tu prendi il secchio e un coltello, e io prendo la batteria e gli altri due coltelli. Penso che questo accordo ci permetta di massimizzare i nostri punti e di avere una buona base per affrontare le sfide future.} \\
\\ 
\texttt{Vorrei fare una proposta segreta per formalizzare questo accordo: {[}Proposta: 1 secchio, 1 coltello, 0 batteria, 0 coltello{]} no, aspetto la tua conferma, poi farò la proposta segreta.} \\
            }
        }
    }
    & & \\ \\

    \theutterance \stepcounter{utterance}  
    & & & \multicolumn{2}{p{0.3\linewidth}}{
        \cellcolor[rgb]{0.9,0.9,0.9}{
            \makecell[{{p{\linewidth}}}]{
                \texttt{\tiny{[GM$|$GM]}}
                \texttt{proposal must include every item type only once} \\
            }
        }
    }
    & & \\ \\

\end{supertabular}
}

\end{document}
